\chapter{Global Specifications}

Global values are controlled in \texttt{config/project\_specs.json}. Do not duplicate payload, dimensions, speed, endurance, or interface values manually in section prose.

\section{Specification Status}

\begin{center}
\fbox{\parbox{0.88\textwidth}{
\textbf{Important specification status}\\

The values in this section are \textbf{provisional study assumptions}. They are not final confirmed product specifications.

The current payload value of \SpecBaseRobotRatedPayloadKg~kg is used only as a starting point for requirements, mechanical sizing, component comparison, safety thinking, and verification planning. The final target payload may remain \SpecBaseRobotRatedPayloadKg~kg, decrease, or increase after the market study, requirements review, mechanical design, electrical validation, cost analysis, and supplier research are completed.
}}
\end{center}

% REVIEW: Do not treat the current payload, dimensions, speed, endurance, or module interface values as final until the project lead approves the final V1 specification baseline.

\section{Base Robot Study Assumptions}

% !TEX root = ../../main.tex
% Generated by tools/generate_specs.py. Do not edit manually.
\small
\begin{longtable}{>{\raggedright\arraybackslash}p{0.24\textwidth}>{\raggedright\arraybackslash}p{0.36\textwidth}>{\raggedright\arraybackslash}p{0.28\textwidth}}
\toprule
Section & Field & Value \\
\midrule
\endhead
project & name & AMR-\allowbreak{}X \\
project & full\_\allowbreak{}name & Autonomous Modular Robot -\allowbreak{} Extended \\
project & version & V1 \\
project & team & AMR-\allowbreak{}X Engineering Group \\
study\_\allowbreak{}status & specification\_\allowbreak{}status & Provisional study assumptions, not final confirmed specifications \\
study\_\allowbreak{}status & payload\_\allowbreak{}status & 50 kg is a starting study assumption only \\
study\_\allowbreak{}status & component\_\allowbreak{}reference\_\allowbreak{}payload\_\allowbreak{}range & 20 to 30 kg \\
study\_\allowbreak{}status & final\_\allowbreak{}target\_\allowbreak{}status & TBD after market, requirements, mechanical, electrical, and cost validation \\
base\_\allowbreak{}robot & rated\_\allowbreak{}payload\_\allowbreak{}kg & 50 \\
base\_\allowbreak{}robot & structural\_\allowbreak{}design\_\allowbreak{}margin & 1.5 \\
base\_\allowbreak{}robot & max\_\allowbreak{}test\_\allowbreak{}load\_\allowbreak{}kg & 75 \\
base\_\allowbreak{}robot & max\_\allowbreak{}length\_\allowbreak{}mm & 800 \\
base\_\allowbreak{}robot & max\_\allowbreak{}width\_\allowbreak{}mm & 600 \\
base\_\allowbreak{}robot & max\_\allowbreak{}base\_\allowbreak{}height\_\allowbreak{}mm & 450 \\
base\_\allowbreak{}robot & min\_\allowbreak{}ground\_\allowbreak{}clearance\_\allowbreak{}mm & 30 \\
base\_\allowbreak{}robot & max\_\allowbreak{}unladen\_\allowbreak{}mass\_\allowbreak{}kg & 80 \\
base\_\allowbreak{}robot & operating\_\allowbreak{}speed\_\allowbreak{}min\_\allowbreak{}mps & 0.5 \\
base\_\allowbreak{}robot & operating\_\allowbreak{}speed\_\allowbreak{}max\_\allowbreak{}mps & 1.5 \\
base\_\allowbreak{}robot & safety\_\allowbreak{}reduced\_\allowbreak{}speed\_\allowbreak{}mps & 0.3 \\
base\_\allowbreak{}robot & endurance\_\allowbreak{}hours & 3 \\
base\_\allowbreak{}robot & min\_\allowbreak{}aisle\_\allowbreak{}width\_\allowbreak{}m & 1.5 \\
base\_\allowbreak{}robot & min\_\allowbreak{}door\_\allowbreak{}width\_\allowbreak{}m & 0.9 \\
base\_\allowbreak{}robot & max\_\allowbreak{}ramp\_\allowbreak{}deg & 5 \\
base\_\allowbreak{}robot & temperature\_\allowbreak{}min\_\allowbreak{}c & 5 \\
base\_\allowbreak{}robot & temperature\_\allowbreak{}max\_\allowbreak{}c & 40 \\
navigation & route\_\allowbreak{}test\_\allowbreak{}distance\_\allowbreak{}m & 50 \\
navigation & waypoint\_\allowbreak{}position\_\allowbreak{}tolerance\_\allowbreak{}mm & 100 \\
navigation & waypoint\_\allowbreak{}heading\_\allowbreak{}tolerance\_\allowbreak{}deg & 5 \\
safety & estop\_\allowbreak{}stop\_\allowbreak{}time\_\allowbreak{}s & 0.5 \\
safety & obstacle\_\allowbreak{}test\_\allowbreak{}count & 20 \\
module\_\allowbreak{}interface & mounting\_\allowbreak{}pattern & TBD \\
module\_\allowbreak{}interface & mounting\_\allowbreak{}grid\_\allowbreak{}spacing\_\allowbreak{}mm & TBD \\
module\_\allowbreak{}interface & max\_\allowbreak{}module\_\allowbreak{}mass\_\allowbreak{}kg & TBD \\
module\_\allowbreak{}interface & power\_\allowbreak{}output\_\allowbreak{}voltage\_\allowbreak{}v & TBD \\
module\_\allowbreak{}interface & power\_\allowbreak{}output\_\allowbreak{}current\_\allowbreak{}a & TBD \\
module\_\allowbreak{}interface & data\_\allowbreak{}interface & TBD \\
module\_\allowbreak{}interface & communication\_\allowbreak{}protocol & TBD \\
\bottomrule
\end{longtable}
\normalsize


\section{Module Interface Specifications}

TBD.

% REVIEW: Define the common mechanical mounting pattern, power output, data interface, communication protocol, safety interface, and maximum allowed module mass.
